\section{Discussion}

HeatGeo is designed around the observation that embedding distillation should match the object that retrieval systems actually use: corpus-level neighborhood geometry. The method has three practical advantages. First, it uses cached teacher embeddings rather than teacher hidden states, attention maps, or online teacher forward passes, which lowers training-time memory pressure. Second, it is tokenizer-agnostic because the loss is defined over final embeddings and graph diffusion targets. Third, it does not change the student architecture, so the distilled model can be deployed as an ordinary text encoder.

The main limitation is that HeatGeo requires a corpus-level preprocessing step. Building a teacher graph is worthwhile when the same training corpus supports substantial distillation, but it adds overhead compared with purely mini-batch objectives. HeatGeo also depends on the quality and coverage of the teacher graph: if the teacher embeddings contain spurious neighborhoods, diffusion may propagate those errors. Future work can study approximate nearest-neighbor construction, adaptive diffusion scales, and online graph refresh for changing corpora.
